\section{Reference case}
\label{sec:reference}

\subsection{Inputs}
The reference case uses the project data already carried on the Losses tab: a
200\,mi line on Bluebird ACSR, two subconductors per pole, with the DMR
provisionally taken as a single Bluebird conductor. Station losses are 0.8\,\%
per station and the current margin 10\,\%.

\begin{table}[htbp]
  \centering
  \renewcommand{\arraystretch}{1.25}
  \begin{tabularx}{\linewidth}{@{}X l r r@{}}
    \toprule
    \textbf{Input} & \textbf{Unit} &
    \textbf{Sym.\ monopole} & \textbf{Bipole} \\
    \midrule
    AC power at receiving end          & MW                & 1{,}500  & 3{,}000  \\
    DC voltage pole-to-ground          & kV                & 400      & 525      \\
    DC line length                     & mi                & 200      & 200      \\
    HV subconductors per pole          & ---               & 2        & 2        \\
    Specific HV resistance             & $\Omega$/1000\,ft & 0.00801  & 0.00801  \\
    DMR subconductors                  & ---               & n/a      & 1        \\
    Specific DMR resistance            & $\Omega$/1000\,ft & n/a      & 0.00801  \\
    Losses per converter station       & \%                & 0.8      & 0.8      \\
    Margin for DC current              & \%                & 10       & 10       \\
    \bottomrule
  \end{tabularx}
  \caption{Reference case inputs.}
  \label{tab:refinputs}
\end{table}

\subsection{Line resistance}
From Equation~\eqref{eq:resistance}, identical for both configurations because
they share the conductor and the length:
\[
  R_{HV} = \frac{0.00801 \times 5.28}{2}\times 200 = 4.229\ \Omega,
  \qquad
  R_{DMR} = \frac{0.00801 \times 5.28}{1}\times 200 = 8.459\ \Omega,
\]
so $R_{loop} = 8.459\,\Omega$ with two poles energised and
$12.688\,\Omega$ with the return through the DMR.

\subsection{Worked step-through, bipole}
Taking the bipole column and following Section~\ref{sec:solution}:
\begin{align*}
  P_{dc,inv} &= \frac{3000}{1-0.008} = 3024.2\ \mathrm{MW}
              && \text{Eq.~\eqref{eq:step1}} \\
  I_{d}      &= \frac{525 - \sqrt{525^{2} - 8.459\times 3024.2}}{8.459}
              = \frac{525 - 500.04}{8.459} = 2.950\ \mathrm{kA}
              && \text{Eq.~\eqref{eq:current}} \\
  P_{dc,rec} &= 2\times 525 \times 2.950 = 3097.8\ \mathrm{MW}
              && \text{Eq.~\eqref{eq:c1}} \\
  P_{loss,line} &= 8.459 \times 2.950^{2} = 73.6\ \mathrm{MW}
              && \text{Eq.~\eqref{eq:c2}} \\
  V_{d,inv}  &= 525 - 2.950\times 4.229 = 512.5\ \mathrm{kV}
              && \text{Eq.~\eqref{eq:vinv}} \\
  P_{ac,rec} &= \frac{3097.8}{1-0.008} = 3122.8\ \mathrm{MW}
              && \text{Eq.~\eqref{eq:c4}} \\
  \eta       &= \frac{3000}{3122.8} = 96.07\,\%
              && \text{Eq.~\eqref{eq:c6}}
\end{align*}
The closure check of Equation~\eqref{eq:closure} balances:
$24.98 + 73.63 + 24.19 = 122.8\ \mathrm{MW} = P_{ac,rec}-P_{ac,inv}$.

\subsection{Results}
\begin{table}[htbp]
  \centering
  \renewcommand{\arraystretch}{1.25}
  \begin{tabularx}{\linewidth}{@{}X l r r r@{}}
    \toprule
    \textbf{Result} & \textbf{Unit} &
    \textbf{Sym.\ monopole} & \textbf{Bipole} &
    \textbf{Bipole, 1 pole out} \\
    \midrule
    Loop resistance                  & $\Omega$ & 8.459   & 8.459   & 12.688  \\
    DC power at inverter terminals   & MW       & 1{,}512.1 & 3{,}024.2 & 1{,}438.5 \\
    DC current, operating            & A        & 1{,}929 & 2{,}950 & 2{,}950 \\
    DC current, rated (10\,\% margin)& A        & 2{,}122 & 3{,}245 & 3{,}245 \\
    DC power at rectifier terminals  & MW       & 1{,}543.6 & 3{,}097.8 & 1{,}548.9 \\
    Line losses, HV conductors       & MW       & 31.5    & 73.6    & 36.8    \\
    Line losses, DMR                 & MW       & ---     & 0.0     & 73.6    \\
    Total DC line losses             & MW       & 31.5    & 73.6    & 110.4   \\
    DC voltage at inverter           & kV       & 391.8   & 512.5   & 487.6   \\
    AC power at receiving end        & MW       & 1{,}500.0 & 3{,}000.0 & 1{,}427.0 \\
    AC power at sending end          & MW       & 1{,}556.0 & 3{,}122.8 & 1{,}561.4 \\
    Total losses, AC to AC           & MW       & 56.0    & 122.8   & 134.4   \\
    Overall efficiency               & \%       & 96.40   & 96.07   & 91.39   \\
    Maximum transferable power       & MW       & 18{,}916 & 32{,}585 & 5{,}431 \\
    Power vs.\ bipole rating         & \%       & ---     & 100.0   & 47.6    \\
    \bottomrule
  \end{tabularx}
  \caption{Reference case results. The one-pole-out column runs at the same
           pole current as the balanced bipole.}
  \label{tab:refresults}
\end{table}

\subsection{Reading the results}
Three points are worth drawing out.

\textbf{The DMR dominates the outage case.} At 110.4\,MW the line loss in the
one-pole-out condition is half again the loss of the balanced bipole at twice
the power, and two thirds of it is in the DMR. That is a direct consequence of
the placeholder assumption of a single Bluebird conductor: the return path is
twice the resistance of the two-conductor bundle it stands in for. Sizing the
DMR closer to the pole bundle moves the residual capability up toward 50\,\%
and cuts the outage loss substantially --- this is the first input to revisit
with real conductor data.

\textbf{Defining power at the receiving end raises the current.} The Losses tab
derives the bipole current directly as $P/(2V_{d}) = 2857\,\mathrm{A}$; the
load flow gives 2950\,A, some 3\,\% higher. The difference is not an
inconsistency but a difference in what the power figure means: the Losses tab
takes 3000\,MW as the power \emph{leaving} the rectifier, whereas here 3000\,MW
must \emph{arrive} at the receiving AC terminals, so the current has to cover
the line and converter losses as well. The load flow figure is the one to size
equipment against when the guarantee is written at the point of delivery.

\textbf{The line, not the converters, sets the loss.} At 0.8\,\% per station
the two converters together account for 49\,MW of the 123\,MW lost in the
bipole case, against 74\,MW in the conductors. Line loss scales with the
square of the current and therefore inversely with the square of the voltage;
converter loss scales linearly with power. That asymmetry is the mechanism
behind the voltage-level trade study, and this tab is the tool for it --- the
comparison between the 400\,kV monopole and the 525\,kV bipole columns is a
first instance of exactly that.
